% =============================================================================
\chapter{Data}
\label{ch:data}
% =============================================================================

% TODO: Structure
%
% \section{Sample Selection}
% - Universe: top ~1,000 U.S. companies by market cap from SEC EDGAR
% - Filters: operating entities only, exclude SIC 6000--6999 (financials)
% - Final sample: ~540 non-financial firms, 2010--2024
% - Table~\ref{tab:sample}: industry distribution, FYE month distribution
%
% \section{Data Sources}
% \subsection{SEC EDGAR}
% - XBRL Company Facts API for financial statements (10-K)
% - Full-text filings for textual analysis (Item 7 MD\&A, Item 1A Risk Factors)
%
% \subsection{Stock Market Data}
% - Source: yfinance (adjusted close prices)
% - Volatility: annualised std of daily log returns over 365-day post-filing window
%
% \section{Financial Variables}
% - Raw XBRL metrics: total\_assets, net\_income, operating\_income,
%   operating\_cash\_flow, eps\_diluted, etc.
% - Derived ratios: log\_total\_assets, leverage, ROA, asset\_growth,
%   current\_ratio, ocf\_to\_assets, operating\_roa
% - Winsorisation at 1\%/99\%
%
% \section{Textual Variables}
% - $\Delta$Sentiment: YoY change in Loughran--McDonald net tone (Item 7)
% - $\Delta$Risk: YoY change in risk-word frequency (Item 1A)
% - TextSim: cosine similarity of consecutive filings
%
% \section{Divergence Measure Construction}
% - Financial trajectory: $\Delta$ROA = ROA$_t$ $-$ ROA$_{t-1}$
% - Narrative trajectory: $\Delta$Sentiment = Sent$_t$ $-$ Sent$_{t-1}$
% - Continuous: Divergence = $\Delta$Sentiment $-$ $\lambda \cdot \Delta$ROA
% - Binary: DivDummy = 1 if sign($\Delta$Sentiment) $\neq$ sign($\Delta$ROA)
%
% \section{Dependent Variable}
% - Primary: annualised volatility (std of daily log returns $\times \sqrt{252}$)
%   over 365-day post-filing window
% - Robustness: post-filing log stock return
%
% \section{Panel Data Construction}
% - Annual panel: ~6,400 firm-years
% - Unbalanced panel, fiscal year-end alignment
%
% \section{Descriptive Statistics}
% - Summary statistics (Table~\ref{tab:desc_stats})
% - Correlation matrix (Table~\ref{tab:correlations})


% ─────────────────────────────────────────────────────────────────────────────
\section{Selecting the Sample}
\label{sec:sample_selection}
% ─────────────────────────────────────────────────────────────────────────────

We start with the sample of 1{,}000 largest U.S.\ public companies by market
capitalisation, taken from the SEC EDGAR company-tickers file, and
apply three filters.  First, we keep only entities the SEC classifies
as ``operating,'' which removes exchange-traded funds, closed-end funds,
blank-cheque companies, and shell entities.  Second, we exclude firms 
in SIC codes 6000–6999 (finance, insurance, and real estate), following 
common sample-selection practice in prior accounting research \parencite{sloan1996,francis2005}. 
Third, we require at least five years of readable XBRL data from
10-K filings and 90\% valid stock-return coverage, so that there is
enough within-firm variation for fixed-effects estimation.

\paragraph{Sample size.}
We deliberately focus on large, liquid firms rather than the full
EDGAR universe for several reasons.  First, our regressions use firm
fixed effects with standard errors clustered by firm.  As
\textcite{petersen2009} shows, the reliability of cluster-robust
standard errors depends on the number of clusters---here 537---not on
total observations, and 537~clusters are well above the
threshold for reliable inference.  Also, the textual-analysis part
of the thesis requires downloading, parsing, and scoring every 10-K in
the sample.  Expanding to several thousand firms would greatly
increase computation time.  Lastly, smaller firms tend to have shorter XBRL histories
(mandatory electronic filing was phased in by firm size), more tagging
errors, and thinner trading, all of which add noise to financial
ratios and return-based variables.

The sample size is in line with prior work combining textual and
financial analysis of annual filings.
\textcite{campbell2014} study about 9{,}000 firm-years of risk-factor
disclosures from large U.S.\ firms; \textcite{loughran2011} use a
broader set of roughly 50{,}000 firm-years, but their goal---validating
a sentiment dictionary---calls for maximal cross-sectional coverage,
whereas our divergence measure is estimated within firms over time.


% ─────────────────────────────────────────────────────────────────────────────
\section{Choosing the Filing Frequency --- why annual 10-Ks?}
\label{sec:annual_vs_quarterly}
% ─────────────────────────────────────────────────────────────────────────────

Before moving to variable construction, it is worth explaining why this
study uses annual 10-K filings instead of quarterly 10-Q reports.
Quarterly data would roughly triple the sample, but two practical
problems---shown in Tables~\ref{tab:q4_reliability}--\ref{tab:text_availability}---ruled
out that approach.

\paragraph{Financial data gaps in Q4.}
10-Q filings cover only Q1--Q3, so any quarterly panel must reconstruct
Q4 as the difference between the annual figure and the sum of the three
preceding quarters.  Table~\ref{tab:q4_reliability} shows that this
residual approach is not usable for operating cash flow: the three
quarterly facts coexist in only $7\%$ of firm-years, because most firms
report cumulative year-to-date cash flows rather than standalone
quarterly amounts.  Operating cash flow is the primary input into our
OCF-to-assets ratio, so a quarterly panel would lose the cash-flow-based
robustness check at the source.  Even for net income and operating
income, where coverage exceeds $80\%$, the implied Q4 disagrees with the
explicit Q4 by more than one percent of the annual total in about
$10$--$14\%$ of comparable firm-years.

\begin{table}[htbp]
\centering
\caption{Q4 Reliability: Availability and Consistency of Fourth-Quarter XBRL Facts}
\label{tab:q4_reliability}
\small
\setlength{\tabcolsep}{5pt}
\begin{tabular}{lrrrrr}
\toprule
Metric & Firm-Years & Q1\,--\,Q3 (\%) & Explicit Q4 (\%) & Comparable & Match (\%) \\
\midrule
Net Income & 451 & 80.9 & 43.0 & 172 & 89.5 \\
Operating Income & 391 & 79.0 & 22.0 & 72 & 86.1 \\
Operating Cash Flow & 462 & 7.4 & 3.2 & 0 & --- \\
\bottomrule
\end{tabular}
\begin{tablenotes}
\small
\item \textit{Note:} ``Q1\,--\,Q3'' indicates the share of fiscal years where all three
quarterly facts (Q1, Q2, Q3) exist in XBRL. ``Explicit Q4'' indicates the share with
a standalone fourth-quarter fact (period $\approx$90~days ending at fiscal year-end).
``Comparable'' is the count of firm-years for which both an explicit Q4 fact and the
implied residual are available. ``Match'' reports the share of those comparable
observations where the two values agree within 1\% of the annual figure.
Sample: 30 non-financial firms, all available fiscal years.
\end{tablenotes}
\end{table}


\paragraph{Missing risk-factor text in 10-Q.}
SEC Regulation S-K only requires firms to update risk factors in
quarterly filings when ``material changes'' have occurred since the
last 10-K.  In practice, only about half the 10-Q filings in our pilot
sample (15~firms, 90~filings) contained a legitimate risk-factor section.
The rest had only a reference back to the annual report
(Table~\ref{tab:text_availability}).  This held across sectors---technology,
healthcare, consumer goods, energy---with few exceptions.  Building
$\Delta$Risk, which enters both H1 and H2, is therefore not feasible
at a quarterly frequency.

\begin{table}[htbp]
\centering
\caption{Textual Data Availability: 10-K vs.\ 10-Q Filings}
\label{tab:text_availability}
\begin{tabular}{lrr}
\toprule
 & 10-K & 10-Q \\
\midrule
Filings tested & 44 & 90 \\
\addlinespace
\multicolumn{3}{l}{\textit{MD\&A (Management's Discussion and Analysis)}} \\
\quad Section found (\%) & 61 & 93 \\
\quad Median word count & 876 & 5,314 \\
\addlinespace
\multicolumn{3}{l}{\textit{Risk Factors}} \\
\quad Substantive section found (\%) & 95 & 36 \\
\quad Boilerplate / reference only (\%) & 0 & 58 \\
\quad Median word count (substantive) & 23,614 & 5,490 \\
\bottomrule
\end{tabular}
\begin{tablenotes}
\small
\item \textit{Note:} ``Found'' requires $>$200 words extracted by regex-based section
parser. ``Substantive'' for 10-Q excludes filings containing only a reference to the
annual 10-K (e.g., ``no material changes from those disclosed in our Annual Report'').
MD\&A corresponds to Item~7 in 10-K and Part~I Item~2 in 10-Q; Risk Factors
corresponds to Item~1A in 10-K and Part~II Item~1A in 10-Q.
Sample: 15 firms, most recent filings.
\end{tablenotes}
\end{table}


\paragraph{Chosen frequency.}
For these reasons---and following prior textual-analysis work based on
10-K filings \parencite{loughran2011,li2008,blankespoor2020}---we use
an annual panel. The XBRL coverage comparison
(Table~\ref{tab:xbrl_coverage}) and a consolidated summary
(Table~\ref{tab:annual_vs_quarterly}) are in
Appendix~\ref{app:diagnostics}.

% Full diagnostic tables deferred to appendix:
% \begin{table}[htbp]
\centering
\caption{XBRL Financial Metric Coverage: Annual (10-K) vs.\ Quarterly (10-Q)}
\label{tab:xbrl_coverage}
\begin{tabular}{lrr}
\toprule
Metric & Annual (\%) & Quarterly (\%) \\
\midrule
assets\_current & 86.7 & 91.1 \\
eps\_diluted & 89.0 & 92.7 \\
liabilities\_current & 86.7 & 91.1 \\
net\_income & 93.2 & 97.3 \\
operating\_cash\_flow & 88.6 & 35.6 \\
operating\_income & 79.7 & 83.8 \\
stockholders\_equity & 98.7 & 94.9 \\
total\_assets & 89.7 & 93.3 \\
total\_liabilities & 66.5 & 69.7 \\
\midrule
Median periods per firm & 20 & 50 \\
Firms with data & 30 & 30 \\
\bottomrule
\end{tabular}
\begin{tablenotes}
\small
\item \textit{Note:} Coverage shows the percentage of firm-periods with non-missing
values for each metric. Annual data is extracted from 10-K filings; quarterly data
from 10-Q filings (Q1--Q3 only). Sample: 30 non-financial operating
companies from the top 200 by market capitalisation.
\end{tablenotes}
\end{table}
       → Appendix
% \begin{table}[htbp]
\centering
\caption{Annual vs.\ Quarterly Approach: Summary Comparison}
\label{tab:annual_vs_quarterly}
\begin{tabular}{p{5.5cm}p{4.5cm}p{4.5cm}}
\toprule
Criterion & Annual (10-K) & Quarterly (10-Q) \\
\midrule
Filing type & 10-K & 10-Q (Q1--Q3 only) \\
MD\&A availability & 93\% (Item 7) & 100\% (Part I, Item 2) \\
Risk Factors availability & 100\% (Item 1A) & $\sim$52\% substantive \\
Risk Factors content & Median $\sim$30{,}000 words & Median $\sim$9{,}000 words (substantive); 48\% boilerplate \\
Q4 financial data & Included in 10-K & Unreliable / absent \\
Explicit Q4 XBRL facts & N/A & 3--43\% of firm-years \\
Implied Q4 mismatch rate & N/A & $\sim$11\% of comparisons \\
Cash-flow quarterly coverage & N/A & 7\% of firm-years \\
NLP variable: $\Delta$Risk & Feasible & Limited (52\% substantive) \\
NLP variable: $\Delta$Sentiment & Feasible & Feasible (noisier) \\
NLP variable: TextSim & YoY (standard) & QoQ (unclear construct) \\
Literature alignment & Standard & Uncommon \\
Expected observations & $\sim$9{,}000 firm-years & $\sim$27{,}000 firm-quarters \\
\bottomrule
\end{tabular}
\begin{tablenotes}
\small
\item \textit{Note:} Summary based on diagnostic tests of 30 non-financial companies
(XBRL, Q4 tests) and 15 companies (textual tests). Risk Factors in 10-Q filings
are required by SEC Regulation S-K Item 1A only when there are ``material changes''
from the most recent 10-K; roughly half of companies include substantive risk
disclosures, while the rest provide only a boilerplate reference. See
Tables~\ref{tab:xbrl_coverage},~\ref{tab:q4_reliability},
and~\ref{tab:text_availability} for detailed results.
\end{tablenotes}
\end{table}
 → Appendix


% ─────────────────────────────────────────────────────────────────────────────
\section{XBRL Tag Consistency Dilemma}
\label{sec:xbrl_tags}
% ─────────────────────────────────────────────────────────────────────────────

A practical problem with the XBRL Company Facts API is that a single
accounting concept can appear under different US-GAAP taxonomy tags.
For example, operating cash flow may be reported as either
\texttt{NetCashProvided\-ByUsed\-InOperating\-Activities} (total, including
discontinued operations) or
\texttt{...ContinuingOperations} (continuing operations only).
Net income may be tagged as \texttt{NetIncomeLoss}, \texttt{ProfitLoss},
or \texttt{IncomeLossFrom\-ContinuingOperations}.
Stockholders' equity can include or exclude the non-controlling
interest.  These tags overlap a lot but are not identical: mixing them all
within a single firm's time series introduces noise,
while picking one tag globally would lose firms that consistently
report under the other.

We handle this problem with a two-part \emph{tag-harmonisation} process:
(i)~\emph{equivalence groups} for tags that are economically
interchangeable, and (ii)~\emph{per-firm tag locking} for tags that
are not.

\paragraph{Equivalence groups.}
When two tags differ only in taxonomy version---not in their economic
meaning---we let both enter a firm's time series.
Specifically, the two operating-cash-flow tags and the pair
\texttt{NetIncomeLoss}/\texttt{ProfitLoss} are treated as
interchangeable.  The third net-income variant,
\texttt{IncomeLossFromContinuingOperations}, is \emph{not} included
because it deliberately excludes discontinued operations and can
differ materially.

To verify this classification, we queried the SEC EDGAR Company~Facts
API for all 550 sample firms and identified 6{,}113 firm-year periods
in which both tags of a pair were reported simultaneously.  The two
overlap samples support treating the pairs as interchangeable:
%
\begin{itemize}
  \item \textbf{Operating cash flow} (1{,}352 overlaps): 53.5\% are
        numerically identical and the median difference scaled by total
        assets is practically zero. The 90th-percentile bias is
        only 1.5~pp.
  \item \textbf{Net income} (4{,}761 overlaps): 25.5\% are numerically identical and
        the median ROA bias is 0.033~pp. 92.5\% of observations lie
        below 1.0~pp.
\end{itemize}
%
The few cases with meaningful differences are
already handled by our priority rule, which prefers the parent-company
tag whenever both are reported.

\paragraph{Per-firm tag locking.}
For metrics where the alternative tags \emph{are} economically distinct
(stockholders' equity with vs.\ without NCI;
\texttt{Income\-Loss\-From\-Continuing\-Operations} vs.\
\texttt{Net\-Income\-Loss}), we apply strict
per-firm tag locking: we keep the tag each firm uses in the most fiscal
years and set values from other tags to missing.  Even if a firm
switches its tag mid-sample, the locked series stays internally
consistent at the cost of some discarded observations.  This
within-firm consistency is what our fixed-effects regressions require.

Combining equivalence groups and per-firm locking increases the
complete-case rate for the five regression controls (\textit{roa},
\textit{leverage}, \textit{asset\_growth},
\textit{current\_ratio}, \textit{log\_total\_assets}) from 78.9\% to
92.7\% of firm-year observations.  This recovers 988~operating-cash-flow
values and 292~net-income values that would otherwise be lost.

\paragraph{Residual limitation.}
The two-part process does not eliminate every source of cross-tag noise.
About 4\% of firms end up with a continuing-operations OCF tag paired
with a total net-income tag.  Ratios that combine the two---such as
the auxiliary OCF-to-assets variable used in robustness--- carry a small inconsistency.  
We accept this trade-off in exchange for better coverage.  
We winsorise these ratios at the 1st and 99th
percentiles to clip the resulting extreme values, and note that the
same kind of cross-tag mismatch is reported in comparable large-scale
XBRL studies \parencite{chychyla2015}.

Per-firm tag choices and the overlap diagnostic are reported in
Appendix~\ref{app:variables}
(Table~\ref{tab:tag_provenance}).


% ─────────────────────────────────────────────────────────────────────────────
\section{How do we extract text from 10-K filings?}
\label{sec:text_extraction}
% ─────────────────────────────────────────────────────────────────────────────

This study heavily depends on the textual sections of 10-K reports Item~1A \emph{Risk Factors} and Item~7
\emph{Management's Discussion and Analysis of Financial Condition and
Results of Operations} (MD\&A).  This section describes how we find
and extract those sections from the raw HTML filings on SEC EDGAR.

\subsection{Source of textual data}
\label{ssec:text_source}

All filings come directly from the SEC EDGAR via the
Submissions~API, which links each company's Central Index Key (CIK)
to the primary documents of its 10-K filings.  From the panel of
about 6{,}700 firm-year observations, we drop 10-K/A amendments.
These typically consist of a cover page and a restated
exhibit and do not contain the full narrative sections of the
original 10-K making them unusable for further analysis.  
Financial data for amended fiscal years are kept,
since the XBRL pipeline already picks up the amended figures when
available (Section~\ref{sec:xbrl_tags}).  Our extraction thus
targets only original 10-K filings.

\subsection{Section Identification}
\label{ssec:section_id}

Unlike structured XBRL data, the textual parts of 10-K filings are hidden
inside HTML with no standardised tagging.  A raw filing can span across
several thousands of lines, with section headers defined in many different
ways: bold text, explicit \texttt{<font>} tags, inline uppercase text,
or table-based layouts. Additionally, these headers are often mixed in with a 
table of contents (TOC), running headers, and inline cross-references (e.g.\ ``see Item~1A'').
These features make simple keyword search on plain text basically impossible,
because it confuses TOC entries, page headers, and inline mentions
with actual section boundaries.

To achieve the best possible results, we use an HTML extraction procedure with four stages:

\begin{enumerate}
  \item \textbf{Header detection.}  We walk through the HTML at
    the block level and test each element against a set of regex
    patterns for the target sections (e.g.\ a pattern that matches
    \emph{``Item 1A''} together with its common variants).
    A candidate for correct section counts only if it meets at least one formatting rule:
    bold styling, ALL-CAPS text, or large font size ($\geq$12\,pt).
    This effectively filters out inline cross-references, which sit in normal
    body-text paragraphs with no special formatting, without discarding any valid headers.  Elements that
    are themselves hyperlinks are also rejected, as these are almost always 
    table-of-contents entries pointing to the real section.

  \item \textbf{Position marking.}  Each confirmed header is tagged
    in place before the document is flattened to plain text, so we
    keep track of where every header sits in the final text.

  \item \textbf{Candidate selection.}  For each target section, we
    keep only headers whose surrounding text contains an expected
    subtitle keyword (``risk'' for Item~1A; ``management'' or
    ``discussion'' for Item~7) and discard any header having many other Item headers near it (within 300~characters in either direction),
    the typical signature of a table of contents.  When several candidates remain, we pick the
    one followed by the longest stretch of text before the next
    section header.

  \item \textbf{Section boundary.}  The section text runs from the
    selected header to the first occurrence of the next expected Item
    header (Item~1B, 1C, or~2 for Item~1A; Item~7A or~8 for Item~7).
    Sections shorter than 200~words are treated as stubs and trigger
    a fallback procedure described in the next subsection.
\end{enumerate}

\subsection{Exhibit Fallback}
\label{ssec:exhibit_fallback}

A common practice in 10-K filings---especially older ones---is to replace an Item with
a short stub such as ``This information is incorporated by reference
from Exhibit~13.''  When the extracted section has fewer than
200~words and the text matches an incorporation-by-reference pattern,
the parser finds the referenced exhibit number (defaulting to
Exhibit~13, the annual report to shareholders) and downloads the
matching HTML document from the filing index.  The section
identification procedure described in the previous subsection is then rerun on the exhibit.  
Because exhibits often skip the formal ``Item~X'' prefix, the parser also
matches patterns with titles only (e.g.\ ``Risk Factors'', ``Management's
Discussion and Analysis'') as a backup detection layer for both the
main document and exhibits.

\subsection{Post-Processing}
\label{ssec:text_postprocessing}

After converting the HTML to plain text, several
layout artefacts from the original filing remain.  These could potentially harm later text analysis.  Before computing any
textual features, we perform 3 cleaning steps
that remove three types of noise:

\begin{enumerate}
  \item \textbf{Running page headers.}  Many filings embed
    ``\textit{$\langle$page number$\rangle$~Table of Contents}''
    strings at page boundaries.  A simple regular expression strips any
    line matching a one- to three-digit number followed by
    ``Table of Contents.''
  \item \textbf{Standalone page numbers.}  After removing the full
    header pattern, some filings still contain page numbers on
    empty lines.  These are also removed 
    by deleting any line consisting solely of a one- to three-digit
    integer.
  \item \textbf{Non-breaking spaces.}  Non-breaking-space characters
    introduced by the original HTML markup are normalised to standard
    spaces, so that downstream word-count and dictionary-based
    measures are not affected by these artefacts.
\end{enumerate}

\noindent
The cleaning rules are intentionally conservative: they target only
patterns that cannot realistically be part of meaningful text and are
applied after section boundaries have already been set, so they do
not interfere with header detection.

\subsection{Coverage and Limitations}
\label{ssec:text_coverage}

Because the quality of the extraction pipeline is crucial for downstream analysis, we validated it on two samples: the 50~most
recent filings in the panel (fiscal year 2024) and, more importantly, a stratified
random sample of 100~filings covering all fiscal years 2010--2024
(20~per era).  On the recent sample, the parser successfully
extracted both sections for all 50~filings (100\%).  On the
cross-era sample, 97~of~100 filings yielded both sections.  The
three failures break down as follows:

\begin{itemize}
  \item Two filings whose annual report exhibit was filed as a PDF
    rather than HTML, making it inaccessible to an HTML parser;
  \item One filing whose 10-K used neither bold, large-font, nor
    uppercase cues for its section headers, leaving no formatting
    signal for the parser to use.
\end{itemize}

\noindent
Neither failure mode is related to firm characteristics in any
systematic way.  On the full panel, we expect a coverage rate of roughly 97\%, with missing values
treated as missing in the textual variable construction.

\paragraph{LLM-assisted quality audit.}
%
% DESIGN RATIONALE (n = 100 filings / 200 sections):
%   1. Statistical precision: With ~200 evaluated sections and an
%      expected pass rate of ~95%, the Wilson 95% CI half-width is
%      z * sqrt(p*(1-p)/n) ≈ 1.96 * sqrt(0.95*0.05/200) ≈ 3.0 pp,
%      yielding intervals narrow enough for a definitive claim.
%   2. Representativeness: 100 filings = 40 sections per 5-era
%      stratum covers the full filing-format evolution
%      (pre-iXBRL → transition → inline XBRL),
%      sampling ~18% of the 540-firm universe.
%   3. Literature precedent: Manual/AI extraction validation in
%      financial NLP typically uses 50–100 documents
%      (cf. Loughran & McDonald spot checks; Krippendorff recommends
%      ≥ 80 units for reliability testing).
%
Although the parser extracted 97/100 filings successfully, coverage
alone says nothing about content quality.  To check whether the
extracted sections are actually the right ones, we ran an independent
validation using a large language model (Google Gemini~2.5~Flash, accessed
through its OpenAI-compatible API, temperature~$= 0$ for deterministic
results).  \textcite{gilardi2023} show that LLM-based annotation
matches or surpasses trained crowd-workers on text-classification
tasks, which supports using a language model as an independent
extraction validator.
A~stratified random sample of 100~filings was drawn---20~filings from
each of five time periods (2010--2012, 2013--2015, 2016--2018,
2019--2021, 2022--2024).  Both Item~1A and Item~7 were evaluated for
every filing, giving 200~section-level evaluations in total.
This design is in line with what \textcite{krippendorff2019} recommends
as reliable for content-analysis studies, and large enough that
pass-rate estimates carry a margin of error of about
$\pm$3~percentage points.  The 100~filings also cover roughly
19\,\% of the 537-firm panel and span the full history of SEC filing
formats.

For each filing, both extracted sections were sent to the model in a
single prompt along with a structured evaluation rubric.  The model
evaluated each extracted section on five criteria---correct section identity,
start-boundary accuracy, end-boundary accuracy, content quality
(no residual HTML tags, page headers, or distorted data), and
completeness (full narrative, not a stub or truncated fragment)---using
a three-level grading scale: \emph{Pass} (criterion fully met),
\emph{Minor} (small issue that would not significantly affect NLP
analysis), and \emph{Fail} (issue that would significantly affect NLP
analysis).  Additionally, we report two aggregate metrics: a strict rate (only Pass) and a lenient
rate (Pass~+~Minor), each with Wilson~95\,\% confidence intervals
\parencite{agresti1998}.

\begin{table}[htbp]
  \centering
  \small
  \caption{LLM-assisted extraction validation results
           ($n = 200$~sections, 100~per section type).}
  \label{tab:ai_validation}
  \begin{tabular}{l r@{/}l r r@{/}l r l}
    \toprule
    Criterion &
      \multicolumn{2}{c}{Strict} & \% &
      \multicolumn{2}{c}{Lenient} & \% &
      95\,\% CI (lenient) \\
    \midrule
    Correct section   & 192 & 200 & 96.0 & 192 & 200 & 96.0 & [92.3\,--\,98.0\,\%] \\
    Start boundary    & 191 & 200 & 95.5 & 193 & 200 & 96.5 & [93.0\,--\,98.3\,\%] \\
    End boundary      & 192 & 200 & 96.0 & 193 & 200 & 96.5 & [93.0\,--\,98.3\,\%] \\
    Content quality   & 191 & 200 & 95.5 & 195 & 200 & 97.5 & [94.3\,--\,98.9\,\%] \\
    Completeness      & 189 & 200 & 94.5 & 191 & 200 & 95.5 & [91.7\,--\,97.6\,\%] \\
    \midrule
    Overall (all five) & 179 & 200 & 89.5 & 187 & 200 & 93.5 & [89.2\,--\,96.2\,\%] \\
    \bottomrule
  \end{tabular}
\end{table}

Table~\ref{tab:ai_validation} shows the results.  On every individual
criterion, the lenient pass rate (Pass~+~Minor) is above 95\%,
confirming that the parser finds the right section with accurate
boundaries and produces clean, complete text.  Even under the strict
metric, the rates stay above 94\% on
each criterion.  The overall lenient rate across all five criteria at
once is 93.5\% (187/200; 95\,\%~CI: 89.2--96.2\,\%).

Of the 13~sections with at least one criteria resulting in Fail, five were complete
extraction failures (empty output) caused by non-standard HTML layouts.
Four had wrong section boundaries (content from an adjacent Item was
included). Finally, the remaining four were flagged only for low word counts
suggesting possible truncation, which we do not permit to enter the main analysis.
Item~7 extractions do slightly better than Item~1A (lenient pass rates
of 96\% vs.\ 91\%), most likely because Item~7 has an advantage in cleaner end
boundary (Item~7A or Item~8), while Item~1A sometimes wrongly picks up text
from the preceding Item~1.

Failures clearly cluster in older filings: 5~of the 40~sections from
2010--2012 were flagged. On the other hand, only 1~of~40 come from 2022--2024.  This
shows expected gradual standardisation of SEC filing HTML over the
sample period.

\paragraph{Robustness of the LLM-based validation.}
Several choices in this validation procedure support the reliability of this approach.
First, the validation is built around clearly defined, verifiable criteria---correct
section identity, boundary accuracy, absence of HTML artefacts, and
word-count thresholds---that leave little room for subjective
interpretation and ambiguity.  Second, each judgment covers a single document with a
deterministic prompt. The model does not need to compare across
filings or aggregate statistics.  Third, manually checking the model's
Fail verdicts and reasons confirmed agreement in every
case examined.

\paragraph{Relation to prior approaches.}
Our extraction strategy is inspired by \textcite{dyer2017}, who
document substantial evolution in 10-K disclosure---length, boilerplate,
and structure---over time and the challenges this creates for
large-sample studies.  Compared with plain-text methods used
in earlier work \parencite{li2008,loughran2011,campbell2014}, our
procedure avoids false positives from TOC entries and inline
cross-references.  Compared with XBRL-tagged textual blocks
(e.g.\ \texttt{us-gaap:RiskFactorsTextBlock}), which became available
with Inline XBRL adoption, our method covers a wider time span:
structured textual tags are inconsistently filled before 2018, while
the HTML-based approach works across the full 2010--2024 window.
